% ======================= COMPILATION NOTICE =======================
% Before compiling, place ge2340_project.sty and all three template
% files (phase1_template.tex, phase2_template.tex, and
% contribution_template.tex) in the same directory.
% Keep the style file named ge2340_project.sty so that the templates
% can load it correctly.
% ==================================================================

\documentclass{article}

% ======================= REQUIRED STYLE =======================
\usepackage[preprint]{ge2340_project}

\usepackage[utf8]{inputenc}
\usepackage[T1]{fontenc}
\usepackage{hyperref}
\usepackage{url}
\usepackage{booktabs}
\usepackage{amsmath,amssymb}
\usepackage{graphicx}
\usepackage{microtype}
\usepackage{xcolor}

% ======================== ASSIGNMENT CONSTRAINTS ========================
% Phase 1 is GROUP WORK. Each group has 3--4 students.
% The proposal must be in English and submitted as a PDF.
% The complete PDF must not exceed 2 pages, including figures and references.
% Do NOT rename, remove, or reorder the required section titles below.
% Replace only the bracketed placeholders and remove instructional text before submission.
% =======================================================================

\title{Phase 1: AI Project Proposal\\
\large Campus Food Chatbot}

% Use the fourth author block only for a four-member group.
\author{
  \textbf{LIU, Cheuk Ying (Student ID: 58538700)}\\
  \texttt{cheuliu2-c@my.cityu.edu.hk}
  \And
  LAU, Wing Lam (Student ID: 58538853)\\
  \texttt{wllau76-c@my.cityu.edu.hk}
  \AND
  WONG, Chung Kai (Student ID: 58529125)\\
  \texttt{kaicwong68-c@my.cityu.edu.hk}
  \And
  LAN, Ka Sen (Student ID: 60166704)\\
  \texttt{fourth.student@my.cityu.edu.hk}
}

\begin{document}
\maketitle

\noindent\textbf{Project topic:} [State the campus or student-life problem in one sentence.]\\
\textbf{Target users/stakeholders:} [Identify who experiences the problem or benefits from the solution.]

% ======================= Required sections (do NOT rename) =======================

\section{Problem Definition \& Scope}
Clearly define a practical problem related to campus or student life. Describe the current situation, affected users, why the problem matters, and the intended scope. State what is included and excluded so that the proposed project is achievable within the course schedule.
What?

[Replace this instruction with your proposal content.]

\section{AI Task Formulation}
Formulate the problem as an AI task. Specify the available or collectable inputs, the expected outputs, the unit of prediction/decision, relevant constraints, and how success will be evaluated. Identify the data source and any privacy, consent, or access considerations.

[Replace this instruction with your proposal content.]

\section{Proposed AI Approaches}
Present multiple plausible AI approaches. For each approach, briefly explain the core method, required data/resources, expected workflow, and practical feasibility. Divergent solutions are encouraged; do not present only one method under different names.

[Replace this instruction with your proposal content.]

\section{Comparative Analysis}
Compare the proposed approaches using relevant dimensions such as expected effectiveness, data needs, implementation effort, compute cost, interpretability, robustness, privacy, and deployment constraints. Clearly identify the preferred approach and justify the choice.

A compact comparison table is encouraged when it saves space.
\begin{table}[t]
  \centering
  \caption{Comparison of candidate AI approaches.}
  \label{tab:approach-comparison}
  \begin{tabular}{lccc}
    \toprule
    Approach & Expected benefit & Main requirement & Main limitation \\
    \midrule
    [A] & [\ldots] & [\ldots] & [\ldots] \\
    [B] & [\ldots] & [\ldots] & [\ldots] \\
    \bottomrule
  \end{tabular}
\end{table}

[Replace this instruction with your proposal content.]

\section{Expected Impact \& Project Feasibility}
Explain the expected practical benefit, novelty, and value to the campus community. Outline the data collection and implementation plan, required tools, major risks, and realistic mitigation measures. Confirm that the group has access to the necessary hardware, software, and data.

[Replace this instruction with your proposal content.]

% ======================= References (included in page limit) =======================
% Cite sources in the text and replace the sample entry below.
\begin{thebibliography}{9}
\small
\bibitem{sample}
[Author(s)], ``[Title],'' \emph{[Venue or Source]}, [Year].

\bibitem{sample}
[Author(s)], ``[Title],'' \emph{[Venue or Source]}, [Year].

\bibitem{sample}
[Author(s)], ``[Title],'' \emph{[Venue or Source]}, [Year].

\end{thebibliography}

% Course policy reminder: LLMs may be used to refine and polish writing,
% but may not be used to generate the submission text from scratch.


\end{document}
